\section{History of Data Models}

% ─────────────────────────────────────────────────────────────────────────────
\begin{frame}{Data Models \& Database Types -- Overview}
  \begin{block}{Central Questions}
    \begin{itemize}
      \item What data models exist, and why were they invented?
      \item What are the conceptual differences between them?
      \item When should you choose one over another?
    \end{itemize}
  \end{block}

  Data models did not appear all at once -- they \emph{evolved} in direct response to
  real-world problems: growing data volumes, the need for concurrent access,
  increasingly complex relationships between entities, and the demand for faster,
  more expressive queries.

  \medskip
  Each model represents a different trade-off between \textbf{simplicity},
  \textbf{flexibility}, and \textbf{scalability}.  Understanding those trade-offs
  is the key skill for any database designer or software engineer:
  the ``best'' model depends entirely on your use case, your data's structure,
  and your query patterns.
\end{frame}

% ─────────────────────────────────────────────────────────────────────────────
\begin{frame}{Paper-Based Data Model}
  \begin{block}{Historical Examples}
    \begin{itemize}
      \item \textbf{Census books} (\textasciitilde{}18th century) -- hand-written population records
      \item \textbf{Marriage registers} (\textasciitilde{}19th century) -- church and civil records of life events
      \item \textbf{Card catalogues} (19th--20th century) -- library index cards, one card per book
    \end{itemize}
  \end{block}

  Paper-based record-keeping served humanity for centuries, but it has severe
  structural limitations.  Finding information required physically searching
  through stacks of books or boxes of cards -- there was no concept of a
  ``query'', no possibility of concurrent access, and no means of automation.

  \begin{alertblock}{Breaking Point: 1890 US Census}
    The 1890 census had to count \textbf{62,947,714} people.  Storing and
    tallying that data manually was no longer feasible -- results from the
    previous census were not fully compiled before the next census began.
    This crisis forced the invention of mechanical data storage and processing.
  \end{alertblock}
\end{frame}

% ─────────────────────────────────────────────────────────────────────────────
\begin{frame}{Herman Hollerith (!)}
  \begin{block}{Mechanical Punch Card Tabulation -- 1890}
    \begin{itemize}
      \item Herman Hollerith invented the \textbf{punch card reader} and
            \textbf{tabulating machine} specifically for the 1890 US census.
      \item Process: an operator feeds punch cards through the reader
            $\rightarrow$ the machine counts holes mechanically
            $\rightarrow$ produces numerical totals.
    \end{itemize}
  \end{block}

  \begin{examples}{}
    Example query: \emph{``How many people in Baltimore are over 50 years old?''}\\
    With Hollerith's machine this could be answered \textbf{mechanically in hours}
    instead of years of manual counting.
  \end{examples}

  Hollerith's invention was transformative not only for data processing but also
  for the computing industry itself.  The company he founded to commercialise
  his tabulating machines eventually merged with others to become the
  \textbf{Computing-Tabulating-Recording Company}, which was renamed
  \textbf{IBM} in 1924.  Punch cards remained a primary data-entry medium
  all the way into the 1970s.

  \smallskip
  \textit{(!) This topic is not part of the exam, but its historical significance
  justifies a brief discussion.}
\end{frame}

% ─────────────────────────────────────────────────────────────────────────────
\begin{frame}{File-Based Data Model}
  \begin{block}{Timeline: 1950s -- 1960s}
    \begin{itemize}
      \item First \textbf{electronic files} stored on magnetic tape and, later, disk
      \item Data is structured by the \emph{application}: a file on a filesystem
            with a format the program understands (fixed-width columns, delimiters, etc.)
      \item Still in everyday use for simple cases: CSV exports, JSON/XML configuration,
            log files
    \end{itemize}
  \end{block}

  The file-based model was a massive leap over paper -- data could now be read,
  written, and processed at electronic speed.  However, every application had
  to implement its own logic for reading and interpreting files.  There was no
  shared schema, no query language, and no access control built in.

  \begin{alertblock}{Key Limitations}
    \begin{itemize}
      \item No concurrent access management -- two programs writing simultaneously
            corrupt the file.
      \item No unified query language -- every developer invents their own search.
      \item Application-specific interpretation -- changing the format breaks
            every program that reads the file.
      \item Data duplication is common because there is no concept of a
            shared, centrally managed data store.
    \end{itemize}
  \end{alertblock}
\end{frame}

% ─────────────────────────────────────────────────────────────────────────────
\begin{frame}{Hierarchical Data Model}
  \begin{block}{Timeline: 1960s -- Developed at IBM}
    \begin{itemize}
      \item Data is arranged in a \textbf{tree structure}: one root node,
            parent nodes, and leaf nodes.
      \item Each record node has $n$ fields; each field has exactly \textbf{1 value}.
      \item Each child has \textbf{exactly 1 parent}; a parent can have $n$ children.
    \end{itemize}
  \end{block}

  \begin{examples}{}
    \texttt{Supermarket} $\rightarrow$ \texttt{Food}, \texttt{Non-Food}\\
    \texttt{Food} $\rightarrow$ \texttt{Cheese}, \texttt{Bread}\\
    \texttt{Cheese} $\rightarrow$ \texttt{Price}, \texttt{Weight}, \ldots
  \end{examples}

  The hierarchical model introduced the concept of \emph{structural relationships}
  between data records.  Navigating from parent to child is very efficient because
  the physical pointers are stored alongside the data.  This made it well-suited
  for \emph{predefined, fixed} traversal paths.

  \begin{alertblock}{Limitation}
    The strict one-parent rule means \textbf{many-to-many relationships cannot be
    represented} naturally.  If a product belongs to two categories, you must
    either duplicate the product record or resort to workarounds -- both
    solutions introduce inconsistency.  The rigidity of the tree structure
    also makes schema evolution painful.
  \end{alertblock}
\end{frame}

% ─────────────────────────────────────────────────────────────────────────────
\begin{frame}{IBM Information Management System (IMS) (!)}
  \begin{block}{Timeline: 1966 by IBM}
    \begin{itemize}
      \item Originally developed to manage the \textbf{Bill of Materials (BOM)}
            for the Apollo moon program -- tracking hundreds of thousands of
            spacecraft components and their sub-components.
      \item Based on the \textbf{hierarchical data model}; processes one dataset
            at a time (no set-based operations).
      \item Programming interfaces: Assembly, C, C++, COBOL, FORTRAN.
    \end{itemize}
  \end{block}

  IMS is one of the oldest database management systems in continuous active use
  -- it has been maintained and extended by IBM for over 50 years.
  Large financial institutions (banks, insurance companies) rely on IMS today
  because their core transaction processing was built on it decades ago and
  the cost of migration exceeds the cost of continued operation.
  This is a powerful real-world lesson: choosing a data model has
  \emph{long-term consequences}.

  \smallskip
  \textit{(!) This topic is not part of the exam.}
\end{frame}

% ─────────────────────────────────────────────────────────────────────────────
\begin{frame}{Network-Like Data Model}
  \begin{block}{Timeline: 1960s -- Charles Bachman at General Electric}
    \begin{itemize}
      \item Data is represented as a \textbf{graph}: \emph{nodes} (records)
            connected by \emph{arcs} (relationships).
      \item Unlike the hierarchical model: a child node can have
            \textbf{multiple parents}.
      \item \textbf{Cycles} in the graph are allowed.
    \end{itemize}
  \end{block}

  \begin{examples}{}
    \texttt{Supermarket} $\leftrightarrow$ \texttt{Food} $\leftrightarrow$
    \texttt{Orders} $\leftrightarrow$ \texttt{Cheese} / \texttt{Bread}\\[2pt]
    A single \texttt{Cheese} node can be linked to both \texttt{Food} and
    \texttt{Gourmet} categories simultaneously.
  \end{examples}

  The network model was a direct solution to the hierarchical model's
  many-to-many limitation.  By allowing a node to have multiple parents,
  it can faithfully represent real-world entities -- for example, a product
  that belongs to several categories, or an employee who reports to
  multiple managers.  The price for this expressiveness is \textbf{navigational
  complexity}: programmers had to write explicit code to traverse pointers,
  which made queries tightly coupled to the physical data structure.
\end{frame}

% ─────────────────────────────────────────────────────────────────────────────
\begin{frame}{Integrated Data Store (IDS) (!)}
  \begin{block}{Timeline: 1964 by General Electric -- Charles Bachman}
    \begin{itemize}
      \item One of the \textbf{first true Database Management Systems (DBMS)}
            in history.
      \item Implemented the network-like data model with physical pointer-based
            navigation between records.
      \item Became the basis for the \textbf{CODASYL} (Conference on Data
            Systems Languages) \textbf{Data Base Task Group} standard -- the
            first attempt to standardise a DBMS interface across vendors.
    \end{itemize}
  \end{block}

  IDS demonstrated that a general-purpose data management layer could exist
  independently of any single application -- a fundamental idea that all
  subsequent database systems inherited.  Charles Bachman received the
  \textbf{Turing Award} in 1973 for his work on IDS and database navigation,
  making him the first database researcher to be so recognised.

  \medskip
  The CODASYL standard eventually influenced the design of COBOL's data
  management facilities, which is why so much legacy financial software
  still bears traces of the network model's thinking.

  \smallskip
  \textit{(!) This topic is not part of the exam.}
\end{frame}

% ─────────────────────────────────────────────────────────────────────────────
\begin{frame}{Relational Data Model}
  \begin{block}{Timeline: 1969 by Edgar F.\ Codd (IBM)}
    \begin{itemize}
      \item Data stored as \textbf{tuples} (rows); tuples grouped into
            \textbf{relations} (tables).
      \item \textbf{Keys} uniquely identify tuples within a relation.
      \item \textbf{Closure property}: the result of any operation on relations
            is itself a relation -- queries can be composed freely.
    \end{itemize}
  \end{block}

  \begin{examples}{}
    \texttt{Person(\underline{PersonID}, FirstName, LastName, Address)}\\
    \texttt{Order(\underline{OrderID}, Date, PersonID\textsuperscript{FK})}
  \end{examples}

  Codd's insight was to \emph{separate the logical view of data from its
  physical storage}.  Applications should describe \emph{what} data they need,
  not \emph{how} to navigate to it.  A high-level, set-based query language
  (eventually SQL) would handle the navigation automatically.
  This \textbf{data independence} made it possible to change storage structures
  without rewriting application code -- a revolutionary idea at the time.

  The relational model became the dominant paradigm through the 1980s and
  remains the most widely used database model today.
\end{frame}

% ─────────────────────────────────────────────────────────────────────────────
\begin{frame}{The Rise of RDBMS in the 1980s}
  \begin{block}{Competing Implementations}
    \begin{itemize}
      \item \textbf{UC Berkeley}: Ingres project $\rightarrow$ eventually became
            \textbf{PostgreSQL} (still one of the most popular open-source RDBMSs)
      \item \textbf{Uppsala University}: research system influencing European DBMS development
      \item \textbf{IBM UK \& IBM San José}: System R project
            $\rightarrow$ gave birth to the \textbf{SEQUEL} query language,
            later renamed \textbf{SQL}
    \end{itemize}
  \end{block}

  Throughout the late 1970s and 1980s, academic labs and IBM research centres
  raced to build working relational systems.  IBM had the research lead but was
  slow to productise -- partly because they feared cannibalising IMS sales.
  \textbf{Oracle} (then called Relational Software, Inc.) recognised the
  commercial opportunity, shipped the first commercial SQL RDBMS in
  \textbf{1979}, and never looked back.

  \begin{alertblock}{Key Lesson}
    Being technically first does not always mean winning the market.  Oracle's
    aggressive go-to-market strategy beat IBM's superior research effort --
    a pattern that repeats throughout technology history.
  \end{alertblock}
\end{frame}

% ─────────────────────────────────────────────────────────────────────────────
\begin{frame}{MySQL}
  \begin{block}{Timeline: 1995 by MySQL AB (now Oracle)}
    \begin{itemize}
      \item Open-source, cross-platform, SQL-based relational database.
      \item \textbf{Fun fact}: ``My'' is the name of co-founder
            Michael Widenius's daughter.
      \item Famous for powering the \textbf{LAMP / WAMP stack}:\\
            Linux (or Windows) + Apache + MySQL + PHP.
    \end{itemize}
  \end{block}

  MySQL's release in 1995 democratised relational databases.  Before MySQL,
  a production-quality SQL database required an expensive commercial licence
  (Oracle, Sybase, Informix).  MySQL was free, fast enough for web workloads,
  and trivial to install -- making it the default choice for the first
  generation of web developers and powering much of the early internet.

  \medskip
  Oracle acquired Sun Microsystems (which had acquired MySQL AB) in 2010.
  Concern about Oracle's stewardship led the original MySQL developers to fork
  the codebase and create \textbf{MariaDB} -- the community-maintained drop-in
  replacement that is now the default in many Linux distributions.
\end{frame}

% ─────────────────────────────────────────────────────────────────────────────
\begin{frame}{XAMPP}
  \begin{block}{Timeline: 2002 by Apache Friends}
    \begin{itemize}
      \item \textbf{X} = Cross-platform (Windows, macOS, Linux)
      \item \textbf{A} = Apache HTTP Server
      \item \textbf{M} = MariaDB (formerly MySQL)
      \item \textbf{P} = PHP
      \item \textbf{P} = Perl
    \end{itemize}
  \end{block}

  XAMPP bundles a complete local web-and-database server stack into a single
  installer.  For learning purposes it is ideal: one download gives you a
  running web server, a relational database, and a scripting runtime without
  any manual configuration.  It ships with \textbf{phpMyAdmin}, a browser-based
  graphical administration tool for MariaDB/MySQL.

  \begin{examples}{}
    In this course we use XAMPP to host a local MariaDB instance.
    You can reach phpMyAdmin at \texttt{http://localhost/phpmyadmin} after
    starting the Apache and MySQL/MariaDB services in the XAMPP Control Panel.
  \end{examples}
\end{frame}

% ─────────────────────────────────────────────────────────────────────────────
\begin{frame}{Object-Oriented Data Model}
  \begin{block}{Timeline: 1980s}
    \begin{itemize}
      \item Maps \textbf{object-oriented programming} concepts directly to
            database storage.
      \item Relation $\equiv$ Class \quad Tuple $\equiv$ Instance \quad
            Column $\equiv$ Attribute \quad Stored Procedure $\equiv$ Method
      \item Gave rise to semi-structured formats such as \textbf{XML} and
            \textbf{JSON} for data exchange and persistence.
    \end{itemize}
  \end{block}

  \begin{examples}{}
    An \texttt{Item} class with attributes \texttt{name}, \texttt{price},
    \texttt{category} can be serialised directly to JSON:\\[2pt]
    \texttt{\{"name":"Gouda","price":3.49,"category":"Cheese"\}}
  \end{examples}

  The object-oriented model emerged because developers found the
  \emph{impedance mismatch} between object-oriented code and relational tables
  painful -- every object had to be manually mapped to rows and columns.
  Object-oriented databases promised to eliminate this friction by persisting
  objects natively.  Adoption was limited because relational systems were
  already entrenched, but the ideas strongly influenced ORM frameworks
  (Hibernate, SQLAlchemy, Django ORM) which bridge the two worlds.
  We will revisit this in the Security \& ORM chapter.
\end{frame}

% ─────────────────────────────────────────────────────────────────────────────
\begin{frame}{ObjectStore (!)}
  \begin{block}{Timeline: 1988 by Versant (originally Progress Software)}
    \begin{itemize}
      \item A database whose unit of storage is a \textbf{C++ object}.
      \item Calling \texttt{new()} in C++ creates a \textbf{persistent object}
            that survives process termination -- no separate ``save'' step.
      \item Use cases: telecommunications, Geographic Information Systems (GIS),
            financial trading systems.
    \end{itemize}
  \end{block}

  ObjectStore showed that the boundary between ``in-memory data structures''
  and ``persistent database records'' could be made nearly invisible to the
  programmer.  The system handled serialisation, indexing, and concurrent
  access automatically.  This concept re-appears in modern object-document
  mappers and in-memory databases like Redis.

  \medskip
  ObjectStore is still in active use in niche domains where the original
  object-oriented codebase is too large and too critical to replace.

  \smallskip
  \textit{(!) This topic is not part of the exam.}
\end{frame}

% ─────────────────────────────────────────────────────────────────────────────
\begin{frame}{1990s -- The Quiet Decade}
  \begin{block}{Consolidation Around the Relational Model}
    \begin{itemize}
      \item No major new database paradigm emerged in the 1990s.
      \item SQL became an ISO/ANSI standard (SQL-89, SQL-92, SQL:1999).
      \item Significant advances in \textbf{data warehousing} and
            \textbf{analytical processing} (OLAP, star schemas, data marts).
      \item MySQL (1995) made relational databases accessible to web developers.
    \end{itemize}
  \end{block}

  The 1990s were important not because of new data models but because of
  \emph{maturation}.  Query optimisers became sophisticated enough that SQL
  performance rivalled hand-tuned navigational code.  Replication, backup,
  and transaction management were standardised.  The internet boom at the end
  of the decade set the stage for the next wave of innovation -- the
  \emph{NoSQL} movement -- which was driven by internet-scale data volumes
  that relational systems struggled to handle.
\end{frame}

% ─────────────────────────────────────────────────────────────────────────────
\begin{frame}{Columnar Data Model}
  \begin{block}{Timeline: 2000s -- Apache Cassandra, HBase, Parquet}
    \begin{itemize}
      \item In a \textbf{row-based} store (traditional RDBMS), each row is
            stored contiguously on disk.
      \item In a \textbf{columnar} store, each \emph{column} is stored
            contiguously -- all values for one attribute sit together.
      \item Reading an analytical query touching only 3 of 100 columns means
            reading $\approx$3\% of the data instead of 100\%.
    \end{itemize}
  \end{block}

  \begin{examples}{}
    \texttt{SELECT AVG(price) FROM products WHERE category = 'Electronics'} \\
    Only the \texttt{price} and \texttt{category} columns need to be read from
    disk -- all other columns are skipped entirely.
  \end{examples}

  Columnar storage also enables \textbf{exceptional compression}: because all
  values in a column share the same data type, techniques like run-length
  encoding and dictionary encoding achieve very high compression ratios.
  The trade-off is that \emph{writing} a single new row requires updating
  many separate column files, making columnar stores slower for
  transactional (OLTP) workloads.  They shine in analytical (OLAP) contexts.
\end{frame}

% ─────────────────────────────────────────────────────────────────────────────
\begin{frame}{Apache Parquet (!)}
  \begin{block}{Timeline: 2013 -- Apache (Twitter \& Cloudera)}
    \begin{itemize}
      \item A \textbf{columnar file format} for big-data processing
            (Apache Spark, Hadoop, Hive, Presto).
      \item Often described as a ``Big Data File Format'': data is spread
            across multiple Parquet files rather than a single monolithic file.
      \item Internal structure:
            \textbf{Header} $\rightarrow$ Row Group 1 $\rightarrow$ Row Group 2
            $\rightarrow \cdots \rightarrow$ \textbf{Footer}
      \item Within each Row Group, column data is stored contiguously,
            enabling efficient compression and selective column reads.
    \end{itemize}
  \end{block}

  Parquet has become the de-facto standard for storing analytical datasets
  in cloud data lakes (AWS S3, Azure Data Lake, Google Cloud Storage).
  Its tight integration with Apache Arrow enables zero-copy reads directly
  into in-memory data frames (pandas, Polars), making the transition from
  storage to computation extremely fast.

  \smallskip
  \textit{(!) This topic is not part of the exam.}
\end{frame}

% ─────────────────────────────────────────────────────────────────────────────
\begin{frame}{Key-Value Data Model}
  \begin{block}{Timeline: Concept from late 1970s; mainstream from 2000s}
    \begin{itemize}
      \item The \textbf{simplest NoSQL model}: every entry is a
            \emph{key} $\mapsto$ \emph{value} pair.
      \item The value can be anything: a string, an integer, a binary blob,
            a serialised object -- the database does not interpret it.
      \item \textbf{No schema}, no relationships, retrieval only by exact key.
    \end{itemize}
  \end{block}

  \begin{examples}{}
    \texttt{"session:u42f8"} $\mapsto$ \texttt{"\{userId:7, cart:[3,11,42]\}"}\\
    \texttt{"rate\_limit:192.168.1.5"} $\mapsto$ \texttt{"47"}
  \end{examples}

  The key-value model's strength is raw speed: lookup by key is an O(1)
  hash-table operation.  This makes it perfect for \textbf{caches},
  \textbf{session stores}, and \textbf{temporary logging}.

  \begin{alertblock}{Key Limitation}
    You can only retrieve data by its exact key.  There is no way to ask
    ``find all sessions belonging to users from Germany'' without scanning
    every entry -- because the database has no knowledge of the value's
    internal structure.  For queries by value, you need a different model.
  \end{alertblock}
\end{frame}

% ─────────────────────────────────────────────────────────────────────────────
\begin{frame}{Redis}
  \begin{block}{Timeline: 2009 by Salvatore Sanfilippo}
    \begin{itemize}
      \item The \textbf{most widely used key-value database} in the world.
      \item \textbf{In-memory} store: all data lives in RAM $\rightarrow$
            sub-millisecond read/write latency.
      \item Optional \textbf{persistence} to disk (RDB snapshots, AOF log).
      \item Supports richer value types beyond plain strings:
            Lists, Sets, Sorted Sets, Hashes, Streams, HyperLogLogs, Bitmaps.
    \end{itemize}
  \end{block}

  Redis is the backbone of many internet-scale architectures.  Common roles:

  \begin{itemize}
    \item \textbf{Cache layer} in front of a relational database
    \item \textbf{Session store} for stateless web applications
    \item \textbf{Pub/Sub message broker} between microservices
    \item \textbf{Rate limiter} (atomic increment on a counter key)
    \item \textbf{Leaderboard} (Sorted Set with score-based ranking)
  \end{itemize}

  A free hosted tier is available at \texttt{redis.io}; the source code is
  open-source on GitHub.
\end{frame}

% ─────────────────────────────────────────────────────────────────────────────
\begin{frame}{Document-Based Data Model}
  \begin{block}{Timeline: Popular from 2000s onward}
    \begin{itemize}
      \item A \textbf{specialised key-value store} where the \emph{value} is a
            structured \textbf{document} (JSON, BSON, or XML).
      \item Documents are grouped into \textbf{collections}; no fixed schema
            -- different documents in the same collection can have different fields.
      \item Unlike plain key-value, you can \textbf{query by any field}
            inside the document.
    \end{itemize}
  \end{block}

  \begin{examples}{}
    Document 1: \texttt{\{name:"Alice", email:"a@x.com"\}}\\
    Document 2: \texttt{\{name:"Bob", email:"b@x.com", phone:"0711...", address:\{...\}\}}\\[4pt]
    Both documents live in the same \texttt{users} collection despite
    having different shapes.
  \end{examples}

  This \textbf{schema flexibility} makes document databases ideal for
  rapidly evolving data (e.g.\ product catalogues where different product
  types have completely different attributes) or for aggregating data from
  many sources.  The trade-off is that consistency enforcement must be
  handled in application code rather than by the database schema.
\end{frame}

% ─────────────────────────────────────────────────────────────────────────────
\begin{frame}{MongoDB}
  \begin{block}{Timeline: 2009 by 10gen (now MongoDB Inc.)}
    \begin{itemize}
      \item Stores documents in \textbf{BSON} (Binary JSON) format --
            a binary-encoded extension of JSON that supports additional types
            (dates, binary data, ObjectId).
      \item \textbf{Hierarchy}: Document $\subset$ Collection $\subset$
            Database $\subset$ Cluster.
      \item Drivers available for all major programming languages
            (Python, Java, Node.js, C\#, Go, PHP, \ldots).
      \item Free tier available at \texttt{mongodb.com} (MongoDB Atlas).
    \end{itemize}
  \end{block}

  MongoDB popularised the document model and the term ``NoSQL'' in the
  mainstream developer community.  It introduced an expressive query language
  (MQL -- MongoDB Query Language) that operates on nested JSON structures,
  aggregation pipelines for complex data transformations, and a flexible
  indexing system.

  \medskip
  Common use cases: content management systems, product catalogues,
  user-profile stores, real-time analytics, and mobile application backends.
\end{frame}

% ─────────────────────────────────────────────────────────────────────────────
\begin{frame}{Graph-Based Data Model}
  \begin{block}{Timeline: 1980s evolution of network model; mainstream 1990s--2000s}
    \begin{itemize}
      \item Data represented as \textbf{Nodes} (entities with properties)
            connected by \textbf{Edges} (relationships with properties).
      \item Both nodes \emph{and} edges can carry arbitrary key-value properties.
      \item Queries traverse the graph to follow relationship chains of
            arbitrary depth.
    \end{itemize}
  \end{block}

  \begin{examples}{}
    Social network: \texttt{(Alice)-[:FRIENDS\_WITH]->(Bob)-[:FRIENDS\_WITH]->(Carol)}\\
    ``Find all friends-of-friends of Alice who live in Stuttgart'' is a
    natural graph traversal, but painful in SQL (recursive CTEs or
    multiple self-joins).
  \end{examples}

  Graph databases excel when \textbf{relationship traversal} is the dominant
  access pattern: social networks, fraud detection (rings of suspicious
  accounts), knowledge graphs, recommendation engines, and network/IT
  infrastructure mapping.  In a relational database, deep traversals require
  expensive recursive joins; in a graph database they are first-class operations.
\end{frame}

% ─────────────────────────────────────────────────────────────────────────────
\begin{frame}{Neo4j}
  \begin{block}{Timeline: 2010 by Neo Technology}
    \begin{itemize}
      \item Query language: \textbf{Cypher} -- a pattern-matching language
            designed to be human-readable (ASCII-art graph patterns).
      \item Components: \textbf{Nodes} (labelled circles), \textbf{Edges}
            (directed, typed relationships), \textbf{Properties}
            (key-value pairs on both).
      \item Industries: insurance, banking, retail, life sciences, cybersecurity.
    \end{itemize}
  \end{block}

  \begin{examples}{}
    \texttt{MATCH (p:Person)-[:BOUGHT]->(i:Item) RETURN p, i}\\[4pt]
    This returns every Person node and the Item nodes they bought, by
    following \texttt{BOUGHT} edges -- equivalent to a SQL JOIN but
    expressed as a graph pattern.
  \end{examples}

  Neo4j is the most widely deployed graph database.  Its native graph storage
  engine stores adjacency lists directly on disk, making multi-hop traversals
  orders of magnitude faster than relational equivalents.  A free ``AuraDB''
  cloud tier is available for experimentation.
\end{frame}

% ─────────────────────────────────────────────────────────────────────────────
\begin{frame}{Summary: Data Models Comparison}
  \begin{block}{Comparison Table}
    \scriptsize
    \begin{tabular}{l l c c c}
      \hline
      \textbf{Data Model} & \textbf{Stored As} & \textbf{Simple} & \textbf{Flexible} & \textbf{Scalable} \\
      \hline
      Paper / File based   & Text                & \checkmark & \checkmark & $\times$ \\
      Hierarchical         & Nodes \& Arcs       & \checkmark & $\times$   & $\times$ \\
      Network-like         & Nodes \& Arcs       & \checkmark & $\times$   & $\times$ \\
      Relational           & Tables              & $\sim$     & $\sim$     & $\sim$   \\
      Object-Oriented      & Objects             & $\sim$     & $\sim$     & $\sim$   \\
      Columnar             & Table / Columns     & $\sim$     & $\times$   & \checkmark \\
      Key-Value            & Key-Value Pairs     & \checkmark & \checkmark & \checkmark \\
      Document             & Documents           & \checkmark & \checkmark & \checkmark \\
      Graph                & Nodes \& Edges      & $\sim$     & \checkmark & \checkmark \\
      \hline
    \end{tabular}
  \end{block}

  \begin{alertblock}{Important Caveat}
    All ratings above are generalisations -- the real answer is always
    \textbf{``it depends on the use case''}.  A relational database can scale
    very well with proper indexing and sharding; a document database can be
    inflexible if you force a rigid schema on it.
  \end{alertblock}

  \begin{block}{Rule of Thumb}
    Choose your data model based on three factors: \textbf{(1)} the natural
    structure of your data, \textbf{(2)} the query patterns your application
    will execute most frequently, and \textbf{(3)} your scalability and
    consistency requirements.  No single model is universally best.
  \end{block}
\end{frame}
